% ════════════════════════ Agradecimientos ═══════════════════════════════
\chapter*{Agradecimientos}
\addcontentsline{toc}{chapter}{Agradecimientos}
\thispagestyle{plain}

\completar{Texto de agradecimientos: tutoría, laboratorio o departamento que ha prestado
el material, familia, compañeros y compañeras.}

\cleardoublepage
% ════════════════════════════ Resumen ═══════════════════════════════════
\chapter*{Resumen}
\addcontentsline{toc}{chapter}{Resumen}
\thispagestyle{plain}

Este Trabajo Fin de Grado aborda el diseño, la construcción y la validación experimental del
sistema electrónico de accionamiento, control y adquisición de un mecanismo de elevación basado
en una \emph{caracola}: un tambor de radio variable que, movido por un servomotor paso a paso de
lazo cerrado, enrolla un cable que levanta una barra articulada con una masa de \SI{2}{\kilo\gram}.
La finalidad de la caracola es modificar la relación de transmisión a lo largo del recorrido para
que el par que debe aportar el motor se mantenga aproximadamente constante, aunque la tensión del
cable varíe con la geometría del mecanismo.

En primer lugar se desarrolla un modelo analítico del mecanismo. Se demuestra que la tensión del
cable es proporcional a la longitud libre de cable, $T = mg(L_3+L_4)\,c/(L_2L_3)$, y se deduce el
perfil de radio que hace el par exactamente constante, $r(\varphi)=C/\sqrt{c_0^2+2C\varphi}$. Se
prueba además que ese par constante, igual al trabajo total dividido por el giro del tambor, es el
mínimo par de pico alcanzable con cualquier perfil. Con la geometría nominal del prototipo, el
perfil ideal reduce el par de pico un \SI{24.7}{\percent} respecto a un tambor cilíndrico que haga
el mismo recorrido, y una espiral de radio lineal optimizada lo reduce un \SI{21.1}{\percent}.

Para verificarlo experimentalmente se ha construido una plataforma completa: una placa adaptadora
diseñada en KiCad alrededor de una Raspberry Pi Pico~2 (bus RS-485, entradas digitales protegidas,
entradas analógicas acondicionadas, salidas de potencia y dos zócalos mikroBUS para células de
carga); un firmware en C++ que gobierna el servo mediante Modbus RTU, lee la célula de carga a
través del acondicionador ZSC31014, realiza el \emph{homing} con dos finales de carrera y ejecuta
ciclos automáticos; un protocolo binario propio con CRC-8; y una aplicación web (FastAPI y
WebSocket) para la supervisión en tiempo real, el registro de ensayos y su comparación con el
modelo.

En los ensayos, el sistema muestrea a \SI{20.0}{\hertz} con una dispersión del periodo de
\SI{0.07}{\milli\second}, alcanza el final de carrera con una desviación típica de
\SI{13}{mrev} y vuelve al origen con un error máximo de \SI{17}{mrev}. Separando
en el par medido la componente gravitatoria y la de pérdidas mediante recorridos en ambos
sentidos, se observa que en el tramo central del recorrido la tensión del cable crece un
\SI{17}{\percent} mientras que el par gravitatorio del motor varía sólo un \SI{6.4}{\percent}, lo
que confirma experimentalmente el efecto compensador de la caracola.

\vspace{1em}
\noindent\textbf{Palabras clave:} tambor de radio variable, caracola, compensación de par,
servomotor paso a paso, Modbus RTU, RS-485, célula de carga, ZSC31014, Raspberry Pi Pico~2,
adquisición de datos, FastAPI.

\cleardoublepage
% ════════════════════════════ Abstract ══════════════════════════════════
\chapter*{Abstract}
\addcontentsline{toc}{chapter}{Abstract}
\thispagestyle{plain}

This Bachelor's Thesis covers the design, construction and experimental validation of the
electronic drive, control and data-acquisition system for a lifting mechanism based on a
\emph{fusee-like spiral drum} (``caracola''): a variable-radius drum, driven by a closed-loop
stepper servomotor, winds a cable that lifts a hinged bar carrying a \SI{2}{\kilo\gram} mass. The
purpose of the spiral drum is to change the transmission ratio along the stroke so that the motor
torque remains approximately constant while the cable tension changes with the geometry.

An analytical model of the mechanism is first derived. The cable tension is shown to be
proportional to the free cable length, $T = mg(L_3+L_4)\,c/(L_2L_3)$, and the radius profile that
makes the torque exactly constant, $r(\varphi)=C/\sqrt{c_0^2+2C\varphi}$, is obtained. This constant
torque, equal to the total work divided by the drum rotation, is proven to be the lowest peak
torque achievable by any profile. With the nominal geometry, the ideal profile lowers the peak
torque by \SI{24.7}{\percent} compared with a cylindrical drum covering the same stroke, and an
optimised linear spiral lowers it by \SI{21.1}{\percent}.

To verify this experimentally, a complete platform was built: a KiCad-designed adapter board
around a Raspberry Pi Pico~2 (RS-485 bus, protected digital inputs, conditioned analogue inputs,
high-side outputs and two mikroBUS sockets for load cells); C++ firmware that commands the servo
over Modbus RTU, reads the load cell through a ZSC31014 conditioner, performs homing with two limit
switches and runs automatic cycles; a custom binary protocol with CRC-8; and a web application
(FastAPI and WebSocket) for real-time monitoring, test logging and comparison with the model.

In the tests, the system samples at \SI{20.0}{\hertz} with a period spread of
\SI{0.07}{\milli\second}, reaches the limit switch with a standard deviation of
\SI{13}{mrev} and returns home with a maximum error of \SI{17}{mrev}. By
separating the gravitational and loss components of the measured torque using strokes in both
directions, it is observed that over the central part of the stroke the cable tension rises by
\SI{17}{\percent} while the gravitational motor torque varies by only \SI{6.4}{\percent},
experimentally confirming the compensating effect of the spiral drum.

\vspace{1em}
\noindent\textbf{Keywords:} variable radius drum, fusee, torque compensation, closed-loop stepper
servo, Modbus RTU, RS-485, load cell, ZSC31014, Raspberry Pi Pico~2, data acquisition, FastAPI.

\cleardoublepage
